% =====================================================================
\chapter{Cuando la IA entra en la formulación: predecir y optimizar}
\label{tema:6}
% =====================================================================

En los temas anteriores los datos eran conocidos. En la práctica la demanda, los tiempos de viaje o los rendimientos se \emph{predicen}, y la decisión hereda el error de la predicción. Este tema estudia cómo combinar ambos pasos. Sigue el capítulo~1 («Introduction and Examples») de Birge y Louveaux~\cite{bl} y las transparencias de programación estocástica del curso EE364a de Stanford~\cite{ee364stoch}; el capítulo~17 («Optimization under Uncertainty») de Kochenderfer y Wheeler~\cite{kw}; los artículos de Elmachtoub y Grigas~\cite{spo} y de Bertsimas y Kallus~\cite{bk} sobre la integración de predicción y decisión; el de Bertsimas y Sim~\cite{bsim} y el libro de Ben-Tal, El Ghaoui y Nemirovski~\cite{btn} para la optimización robusta, y el de Smith y Winkler~\cite{smith} para la maldición del optimizador.

\section{Predecir y después optimizar}

\begin{definicion}[Predecir y después optimizar]\label{def:pto}
Sea un problema $\min_{x\in\mathcal{X}}\E[c(x,\xi)\mid w]$, donde $\xi$ es un parámetro incierto (la demanda) y $w$ el contexto observado (temperatura, festivos). El enfoque \emph{predecir y después optimizar} entrena un modelo $\hat\xi=m_\phi(w)$ y resuelve $\hat x(w)\in\argmin_{x\in\mathcal{X}}c(x,\hat\xi)$, sustituyendo el parámetro por su predicción puntual.
\end{definicion}

\begin{proposicion}[Cuándo basta la media]\label{prop:lineal}
Si el coste es lineal en el parámetro, $c(x,\xi)=\xi\T x$, y $\mathcal{X}$ no depende de $\xi$, la decisión que minimiza el coste esperado condicionado es la que resulta de sustituir $\xi$ por su media condicionada $\E[\xi\mid w]$.
\end{proposicion}

La proposición explica el éxito del enfoque cuando la incertidumbre está en los costes de un PL (rutas con tiempos inciertos y objetivo lineal) y su fracaso cuando el coste es no lineal en $\xi$: con inventario, el coste de quedarse corto y el de sobrar son distintos, y la media deja de ser suficiente. Incluso en el caso lineal, dos modelos con el mismo error cuadrático pueden inducir decisiones de calidad muy distinta: lo que importa no es el error en la predicción, sino el error en la decisión, que es lo que mide la pérdida SPO (\emph{smart predict-then-optimize}) de Elmachtoub y Grigas~\cite{spo} (\cref{eje:6-spo}). Bertsimas y Kallus~\cite{bk} proponen, en lugar de predecir un valor, aprender pesos sobre escenarios históricos parecidos al contexto y optimizar contra esa distribución empírica.

\section{El vendedor de periódicos}

\begin{definicion}[Problema del vendedor de periódicos]\label{def:newsvendor}
Se decide una cantidad $q\ge 0$ antes de conocer la demanda $D$, con función de distribución $F$. Cada unidad que falta cuesta $c_u>0$ (margen perdido) y cada unidad que sobra, $c_o>0$ (coste menos valor de saldo). El coste esperado es
\[
G(q)=\E\bigl[c_u(D-q)^{+}+c_o(q-D)^{+}\bigr].
\]
\end{definicion}

\begin{teorema}[Fractil crítico]\label{teo:fractil}
Si $F$ es continua, $q^{*}=F^{-1}\!\bigl(c_u/(c_u+c_o)\bigr)$ minimiza $G$. Si $D$ es discreta, es óptimo el menor $q$ con $F(q)\ge c_u/(c_u+c_o)$.
\end{teorema}

Si la demanda es normal, $q^{*}=\mu+\sigma\,\Phi^{-1}(c_u/(c_u+c_o))$: se produce por encima de la media cuando el margen perdido supera al coste del sobrante. Predecir la media, por precisa que sea, no basta; hace falta un \emph{cuantil}, y por tanto una estimación de la incertidumbre calibrada fuera de muestra.

\section{Programación estocástica y robusta}

\begin{definicion}[Programa estocástico en dos etapas]\label{def:dos-etapas}
En un \emph{programa en dos etapas} se fija una decisión $x$ antes de observar $\xi$ y, una vez observado, una decisión de recurso $y$:
\[
\mathrm{RP}=\min_{x\in\mathcal{X}}\ c\T x+\E_\xi\bigl[Q(x,\xi)\bigr],\qquad Q(x,\xi)=\min_{y\in\mathcal{Y}(x,\xi)}q(\xi)\T y .
\]
La \emph{aproximación por media muestral} (SAA) sustituye la esperanza por la media sobre $S$ escenarios $\xi^{1},\dots,\xi^{S}$ y, si todo es lineal, da un único PL con una copia de $y$ por escenario~\cite{ee364stoch}.
\end{definicion}

\begin{definicion}[Valor de la solución estocástica y de la información perfecta]\label{def:vss}
Para un problema de minimización, sea $\bar x$ la solución del \emph{problema del valor esperado} (sustituir $\xi$ por $\E\xi$), $\mathrm{EEV}=c\T\bar x+\E[Q(\bar x,\xi)]$ su coste esperado real, y $\mathrm{WS}=\E_\xi[\min_x(c\T x+Q(x,\xi))]$ el coste esperado si se conociera $\xi$ antes de decidir (\emph{wait-and-see}). Se definen el \emph{valor de la solución estocástica} $\mathrm{VSS}=\mathrm{EEV}-\mathrm{RP}$ y el \emph{valor esperado de la información perfecta} $\mathrm{EVPI}=\mathrm{RP}-\mathrm{WS}$~\cite{bl}.
\end{definicion}

\begin{proposicion}\label{prop:vss}
$\mathrm{WS}\le\mathrm{RP}\le\mathrm{EEV}$; en consecuencia, $\mathrm{VSS}\ge 0$ y $\mathrm{EVPI}\ge 0$.
\end{proposicion}

La \emph{optimización robusta} renuncia a la distribución y exige factibilidad o buen coste para todo $\xi$ de un \emph{conjunto de incertidumbre} $\Xi$: $\min_x\max_{\xi\in\Xi}c(x,\xi)$. Para restricciones lineales con coeficientes en una caja, el problema robusto sigue siendo lineal.

\begin{proposicion}[Contrapartida robusta de una caja]\label{prop:robusta}
Sea $x\ge 0$ y $a_j\in[\bar a_j-\delta_j,\ \bar a_j+\delta_j]$ con $\delta_j\ge 0$. La restricción $a\T x\le b$ se cumple para todos los $a$ de la caja si y solo si $(\bar a+\delta)\T x\le b$. Para $x$ de signo libre, la condición es $\bar a\T x+\delta\T\abs{x}\le b$, que se linealiza con la \cref{prop:abs}.
\end{proposicion}

La caja completa es muy conservadora: supone que todos los coeficientes toman su peor valor a la vez. Bertsimas y Sim~\cite{bsim} limitan a $\Gamma$ el número de coeficientes que se desvían y controlan así el \emph{precio de la robustez}: el coste que se paga por protegerse.

\begin{observacion}[La maldición del optimizador]\label{obs:maldicion}
Si se eligen la mejor de varias alternativas por su valor estimado, el valor estimado de la elegida sobrestima su valor real, aunque cada estimación sea insesgada: $\E[\max_i\hat J_i]\ge\max_i\E[\hat J_i]$~\cite{smith} (\cref{eje:6-maldicion}). Por eso el beneficio que promete un optimizador sobre datos predichos debe validarse con datos nuevos, nunca con los usados para elegir.
\end{observacion}

\section{Caso de estudio: plan de cerveza fresca con demanda predicha}

\addtocounter{definicion}{1}%

\begin{table}[h]
\centering\small
\caption{Tecnologías del tema~6: predecir y optimizar bajo incertidumbre.}\label{tab:tec6}
\begin{tabularx}{\textwidth}{@{}l>{\raggedright\arraybackslash}X>{\raggedright\arraybackslash}p{2.2cm}>{\raggedright\arraybackslash}X@{}}
\toprule
Herramienta & Para qué sirve & Licencia o tipo & Cuándo elegirla \\
\midrule
scikit-learn & regresión, cuantiles (\texttt{GradientBoostingRegressor} con pérdida cuantílica), validación temporal & abierta (BSD-3) & modelos de predicción tabulares \\
SciPy / HiGHS & PL de SAA con escenarios & abierta (BSD-3; MIT) & programas estocásticos lineales de tamaño medio \\
CVXPY~\cite{cvxpy} & modelos robustos y con restricciones probabilísticas convexas & abierta (Apache~2.0) & contrapartidas robustas cónicas, CVaR \\
Pyomo & modelos estocásticos por escenarios y descomposición & abierta (BSD-3) & problemas grandes en dos o más etapas \\
PyTorch & capas de optimización diferenciables para aprender con la pérdida de decisión & abierta (BSD) & enfoques SPO o de aprendizaje orientado a la decisión \\
\bottomrule
\end{tabularx}
\end{table}

\addtocounter{definicion}{1}%

